% 第 7 章 Fixing and Optimizing Grid and Power Performance（原书 7-122 … 7-170）
\bichapter{网格与电源性能的修复与优化}{Fixing and Optimizing Grid and Power Performance}
\label{chap:grid}

\section{引言}
\subsection*{Introduction}

RedHawk 静态（IR）或动态压降（DvD）分析识别出设计中不可接受压降区域后，RedHawk Fixing and Optimization（FAO）功能可修改电源网格和/或分配新去耦电容，以满足特定静态与动态压降目标，以及电源相关电路时序问题。修复压降问题亦可减轻地弹噪声。多 Vdd 电源 instance 的操作可透明处理。

在 trial 电源网格与初始布局信息可用后，RedHawk FAO 可全局修改电源网格参数以满足特定压降目标。RedHawk 可在运行任何阶段方便地添加或删除电源/地 pad、strap 与 via，以探索其对压降与电迁移（EM）的影响。RedHawk 可在 IR drop 与 EM 分析过程中增量接纳新变更。详细布局之后、详细布线之前，RedHawk 可识别高动态压降区域或 DvD 对时序影响显著的区域，并通过电源网格导线修改与 targeted decap 布局进行修复。

本章第一部分描述 RedHawk 中可用的手动网格与 decap 修改命令：
\begin{itemize}
  \item 修改现有层或 via 电阻率
  \item 添加电源/地 pad
  \item 删除电源/地 pad
  \item 添加 via
  \item 删除 via
  \item 添加电源 strap
  \item 编辑/删除电源 strap
  \item 添加 decap 单元
  \item 添加金属层与 via 或 via 阵列
  \item 写出 ECO 文件
  \item 读取 ECO 文件
\end{itemize}

设计进入最终布局与布线后，可使用 RedHawk FAO 自动修改网格参数与 decap 布局，以满足压降与电源相关电路时序目标。所有压降约束满足后，RedHawk 可导出含所需设计变更的 ECO 文件，并对设计执行最终 Spice 精度功耗分析 signoff。

FAO 变更后可执行 Undo 或 Redo，便于尝试后撤销不满意的结果。可连续执行任意次数 Undo/Redo，直至命令栈耗尽。

\section{手动电源网格修改}
\subsection*{Manual Power Grid Modification}

修复与优化功能见以下各节。

\subsection{修改金属层或 Via 电阻率}
\subsubsection*{Changing a Metal Layer or Via Resistivity}

可在 RedHawk technology（\texttt{.tech}）文件中修改现有金属层或 via 电阻率，然后重跑 RedHawk 观察影响。Apache \texttt{.tech} 文件格式见附录 C「File Definitions」。

\subsection{添加单个电源/地 Pad}
\subsubsection*{Adding a Single Power/Ground Pad}

使用 Edit $>$ Add Pad 创建新电源或地 pad（图~\ref{fig:rh-add-pad}）。弹出菜单选择 pad 所连金属层，在版图上点击目标位置。新 pad 添加到所选位置。新电源 pad 显示为橙色矩形框，新地 pad 为白色框。使用 File $>$ Export ECO 导出新增 pad，File $>$ Import ECO 导入以回注后续运行。详见本节「使用 ECO 命令保存设计变更」。

\figplaceholder{Figure 7-1}{Edit $>$ Add Pad 菜单}
\label{fig:rh-add-pad}

\subsection{在指定区域添加一组电源/地 Pad}
\subsubsection*{Adding a Set of Power/Ground Pads over a Specified Area}

命令 \cmd{fao add pad} 在整个 mesh 或定义窗口内添加一组电源/地 pad。须指定窗口、金属层、网络，以及 x、y 方向 pitch，在区域内均匀分布 pad。支持 Undo/Redo。命令行语法：
\begin{lstlisting}
fao add pad -window {llx lly urx ury} -metal <layer>
      -net <net_name> -pitchX <microns> -pitchY <microns>
      -r <resistance> -l <inductance> -c <capacitance>
\end{lstlisting}
其中：
\begin{itemize}
  \item \opt{-window \{ \}}：矩形区域左下与右上坐标（默认整个 mesh）
  \item \opt{-metal}：布线层（必需）；每层独立放置
  \item \opt{-net}：网络名（必需）
  \item \opt{-pitchX -pitchY}：自 \opt{-window} 左下角起 pad 间距（必需）
  \item \opt{-r -l -c}：pad RLC 参数（单位与 PLOC 文件相同）
\end{itemize}

\subsection{删除电源/地 Pad}
\subsubsection*{Deleting a Power/Ground Pad}

使用 Edit $>$ Delete Pad 选择并删除现有电源/地 pad。

\subsection{添加 Via}
\subsubsection*{Adding a Via}

使用 Edit $\rightarrow$ Add Via，在版图上点击目标位置。弹出表单定义插入 via stack 的上下金属层。添加叠层 via 时须选择正确数量的 via 与 via 层。新 via 保存后显示为蓝色框。

\subsection{删除 Via}
\subsubsection*{Deleting a Via}

使用 Edit $>$ Delete Via。可能须通过 View Layers Configuration 或 View $\rightarrow$ Map Configuration $\rightarrow$ Layers Color Map 打开 via 显示。左键点击选中 via（黄色高亮），弹出菜单选择删除。

\subsection{添加一条或多条电源 Strap}
\subsubsection*{Adding One or Multiple Power Straps}

\begin{enumerate}
  \item 图形添加（推荐）：Edit $\rightarrow$ Add Power Strap，显示 Power Strap 对话框。
  \item 取消勾选 Add strap by text input。
  \item 显示与新 strap 相关的层。
  \item 右键在版图上绘制 strap。若需多条相同 strap，先绘制代表一组外缘的一条，并勾选 Add multiple straps。
  \item 在对话框中计算并输入最后一条 strap 外缘位置与 pitch。RedHawk 在约束内添加最大合法 strap 数。
  \item 选择添加层（默认 metal1）。
  \item 规格与绘制正确后点击 Commit adding（无 undo）。物理错误或短路将显示错误消息。成功示例如图~\ref{fig:rh-multiple-straps}。
  \item 若需修改，使用 Edit$\rightarrow$Power Strap 选择后删除或修改 strap。
\end{enumerate}

\figplaceholder{Figure 7-2}{按文本输入添加 strap 的表单}

\figplaceholder{Figure 7-3}{使用 Add multiple straps 添加多条 strap}
\label{fig:rh-multiple-straps}

添加 strap 示例见第~\ref{chap:static}~章「修复 IR Drop 问题的示例流程」。

\paragraph{GUI 选项 \cmd{eco add strap}}
勾选 Include stackvias 可在 GUI 上为布线上下两层或更多层自动添加叠层 via。

\figplaceholder{Figure 7-4}{\cmd{eco add strap} 的 GUI 选项}

行为：
\begin{itemize}
  \item 勾选 Include stackvias：\cmd{eco add strap} 为所有可能金属层添加 via。例如在 M2 添加导线 A 时，搜索除 M2 外所有金属层上与 A 重叠的导线并尽可能添加 via。
  \item 未勾选：仅相邻层。例如在 M2 添加导线 A 时，仅搜索 M1/M3 上与 A 重叠的导线。
\end{itemize}

\subsection{编辑/删除电源 Strap}
\subsubsection*{Editing/Deleting a Power Strap}

Edit $>$ Edit Power Strap 选择并编辑现有 strap。左键点击选中（黄色高亮），弹出菜单可修改长度、宽度及 x、y 起点（图~\ref{fig:rh-edit-strap}）。亦可从设计数据库删除 strap。

\begin{noteBox}
编辑现有电源 strap \textbf{不会}自动添加新 via。
\end{noteBox}

\figplaceholder{Figure 7-5}{Edit$\rightarrow$Edit Power Strap}

\subsection{添加 Decap 单元}
\subsubsection*{Adding a Decap Cell}

Edit $>$ Add Decap Cell 添加一个或多个 intentional decap。在表单上点击 Select Vdd，再在最低金属层上点击 VDD 导线位置；点击 Select Gnd 并在邻近位置点击 VSS。表单显示 GSR 中 \gsr{DECAP_CELL} 定义的可用 decap 模型；若无须在 GSR 中提供。

\begin{noteBox}
LEF 与 GSR 中须有对应 DECAP 单元定义。
\end{noteBox}

intentional decap 经 RedHawk APL 流程表征后，从 \texttt{<cell>.cdev} 获取 ESC 与 ESR。点击 Commit Adding 在版图上显示 decap（图~\ref{fig:rh-add-decap}、图~\ref{fig:rh-decap-layout}）。使用 File $\rightarrow$ Export ECO 写出新增 decap。

\figplaceholder{Figure 7-6}{Edit $\rightarrow$ Add Decap Cell}

\figplaceholder{Figure 7-7}{版图上已添加的 decap}

\subsection{添加金属层与 Via 或 Via 阵列}
\subsubsection*{Adding Metal Layers and Vias or Via Arrays}

在现有金属层之上添加 strap 的步骤（例如现有 metal1--metal6，新增层为 metal7）：

\begin{enumerate}
  \item 在 .lef 的 LAYER 段添加层定义，例如添加 VIA67 与 METAL7（顺序与编号须正确）：
\begin{lstlisting}
LAYER VIA67
   TYPE CUT ;
END VIA67
LAYER METAL7
   TYPE ROUTING ;
END METAL7
VIA via6 DEFAULT
RESISTANCE 2.5400e+00 ;
LAYER METAL6 ;
   RECT -0.240 -0.190 0.240 0.190 ;
LAYER VIA67 ;
   RECT -0.180 -0.180 0.180 0.180 ;
LAYER METAL7 ;
   RECT -0.270 -0.270 0.270 0.270 ;
END via6
\end{lstlisting}
  在 \texttt{*.tech} 中添加金属层与 via 信息。
  \item 在顶层 .def 的 VIAS 段创建所需 via 阵列（可复制 VIA56 定义并改为 VIA67 参数）。
  \item 金属层与 via 阵列创建后，按前述步骤添加 pad 与 strap。
\end{enumerate}

\subsection{Undo 与 Redo}
\subsubsection*{Undo and Redo}

FAO 变更后可 Undo/Redo。GUI：Edit $\rightarrow$ Undo / Edit $\rightarrow$ Redo。TCL：\cmd{fao undo}、\cmd{fao redo}。

\section{自动网格优化与修复流程}
\subsection*{Automated Grid Optimization and Fixing Procedures}

\subsection{静态 IR Drop 改进的修复与优化流程}
\subsubsection*{Fixing and Optimization Flow for Static IR Drop Improvement}

FAO 静态压降优化主流程见图~\ref{fig:rh-fao-static-flow}。

\figplaceholder{Figure 7-8}{RedHawk 静态 IR drop 分析流程}
\label{fig:rh-fao-static-flow}

\subsection{多 Vdd 设计的 FAO 流程}
\subsubsection*{FAO Procedure for Multiple Vdd Designs}

多 Vdd 域设计中，FAO 每次仅分析并修改一个电源域。用 GSR \gsr{FAO_NETS} 依次指定各域，例如：
\begin{lstlisting}
gsr set fao_nets VDD
\end{lstlisting}
运行 FAO 并修改 VDD 网格与 decap 后：
\begin{lstlisting}
gsr set fao_nets VDD_B
\end{lstlisting}
再对 VDD\_B 运行 FAO。

若 \gsr{FAO_NETS} 中定义多于 1 个电源网络，仅考虑第一个，并返回警告：
\begin{lstlisting}
WARNING: *** FAO_NETS has > 1 VDD net: Only consider <VDD>
WARNING: *** Set fao_nets for EACH power net and apply fao separately.
\end{lstlisting}

\subsection{网格优化}
\subsubsection*{Grid Optimization}

\paragraph{Mesh 命令}
FAO 调整网格宽度的关键命令（GUI 或 TCL）：
\begin{itemize}
  \item \cmd{mesh optimize}——在 \gsr{fao_region} 内搜索网格宽度解；指定金属层上所有线段调整为相同宽度以解决压降（全芯片或区域内全高度）。\opt{-taper} 时仅在 \gsr{fao_region} 内调整，不横跨全芯片高度。
  \item \cmd{mesh fix}——在包含所有热点及其 fix\_window 的矩形区域内搜索解；热点所在垂直/水平带（全芯片宽/高，由 fix\_window 尺寸定义）上调整导线宽度。\opt{-taper} 时仅在指定区域内调整。
  \item \cmd{mesh snscalc}——计算并报告各层、各网络修改导线宽度的相对平均灵敏度（mv/$\mu$m），数值越高表示加宽导线对降低最差压降越有效。建议在网格缩放前使用。
  \item \cmd{mesh sub_grid}——按 \gsr{fao_region}、\gsr{fao_sub_grid_nets}、\gsr{fao_sub_grid_spec} 添加子网格，目标压降由 \gsr{noise_reduction} 指定；\gsr{fao_range} 指定新网格最小/最大宽度。
  \item \cmd{mesh set_width}——修改指定层 mesh 导线宽度；默认两侧等量扩展，可用 \texttt{expand\_dir} 仅一侧扩展或收缩。
\end{itemize}

其他命令：\cmd{mesh add}、\cmd{mesh delete}、\cmd{mesh vias}、\cmd{mesh generate}、\cmd{ring add/delete}。

GSR 关键字通用 TCL 语法：\cmd{gsr set <variable name> <value>}。详细说明见附录 C「Grid Fixing and Optimization Keywords」。

\paragraph{基于区域的网格宽度调整 GSR 关键字（mesh optimize）}
\begin{enumerate}
  \item \gsr{FAO_DYNAMIC_MODE}（默认 off）——选择 IR 或 DvD 分析
  \item \gsr{NOISE_REDUCTION}（必需）——目标压降降低量
  \item \gsr{FAO_NETS}（必需）——优化网络
  \item \gsr{fao_layers}、\gsr{FAO_RANGE}（均必需）——层与宽度约束
  \item \gsr{FAO_TURBO_MODE}（默认 on）——accurate 或 turbo 模式
  \item \gsr{FAO_REGION}（默认全芯片）——优化区域
  \item \gsr{NUM_HOTSPOT}（默认 1）——优化的高压降区域数
  \item \gsr{FAO_WIDTH_CNSTR}——层/宽度关系
  \item \gsr{MESH_SEARCH}——各层压降降低百分比分配
\end{enumerate}

\paragraph{基于热点的网格宽度修复 GSR 关键字（mesh fix）}
上述 1--10 亦适用于 \cmd{mesh fix}。\gsr{fao_region} 定义总评估区域；每个热点周围的 fix window 内调查网格修改：
\begin{enumerate}
  \setcounter{enumi}{10}
  \item \gsr{fix_window}（默认 100$\mu$m $\times$ 100$\mu$m）——以热点为中心修改网格宽度的区域尺寸
\end{enumerate}

\subsection{网格优化示例}
\subsubsection*{Examples of Grid Optimization}

\paragraph{示例 A——全芯片优化：放宽静态 IR Drop}
条件：设计早期，可接受更高最坏压降以节省网格资源。使用 \cmd{mesh optimize} 全芯片均匀优化。目标：将静态 IR drop 目标提高 15\%。
\begin{enumerate}
  \item GSR 设置（初始最坏静态 IR 18\,mV on VDD）：
\begin{lstlisting}
gsr set fao_turbo_mode 0
gsr set fao_nets VDD
gsr set fao_layers { {metal6 hor <= 29.4} }
gsr set fao_range { {metal6 8 29.4}}
gsr set noise_reduction -15
\end{lstlisting}
  未定义 \gsr{fao_region} 则优化整个电网。FAO 以 accurate 模式搜索低压降区域，仅优化 VDD 上宽度 $\leq$29.4\,$\mu$m 的 METAL6 水平线，尝试将最坏静态压降放宽 15\%（须指定小于现有宽度的范围）。
  \item \cmd{mesh optimize -eco <filename>.eco}
  \item 根据优化参数，RedHawk 在定义区域内报告优化导线，例如：
\begin{lstlisting}
ECO File: fao_mesh_opt_23:15:47_022242005.eco
Region: <0 0 6817.32 8316.08>
FAO Mode: accurate
Simulation Mode: static
Net: VDD
Noise Reduction: -15%
+++++++++++++++++++++++++++++++++++++++
Initial Worst Static Noise: 17.578mv
Voltage Noise Constraints: 20.2147mv
*============================================*
*               Searching Target Grids                         *
*============================================*
Net <VDD>:
     Layer <metal6>: identified 44 wires
*============================================*
*                Sizing Target Grids                           *
*============================================*
Optimization succeed: no error
Optimized width of layer metal6: 23.8um
*========================================================*
 Metal Usage Change Report
      Net <VDD>:
            Layer <metal6>:5.05281e+06um^2 => 4.09119e+06um^2 (-19.03%)
      Total: 5.05281e+06um^2 => 4.09119e+06^2 (-19.03%)
*========================================================*
Replacing optimized grids...
Replace <metal6> mesh grids :
      net <VDD> : 44 wires replaced.
DEF file for new via models: SH_VIA_VDD.def
\end{lstlisting}
  FAO 报告 44 条导线变更建议；放宽静态压降限后 metal6 VDD 金属用量减少，并指定匹配导线变更的新 via 模型。
  \item 最终验证：\cmd{perform extraction}、\cmd{setup pad}、\cmd{setup wirebond}、\cmd{setup package}、\cmd{perform analysis -static}（封装 RLC 单位分别为 $\Omega$、pH、pF）。示例结果：VDD 最终静态压降 21\,mV（放宽 15\%），整片金属资源节省 19\%。
\end{enumerate}

\paragraph{示例 B——全芯片网格修复：降低静态 IR Drop}
条件：需更低压降。目标：自现有 18\,mV 降低静态 IR drop（图~\ref{fig:rh-static-ir-map}）。
\begin{enumerate}
  \item 收紧压降时须增大 \gsr{fao_range}：
\begin{lstlisting}
gsr set fao_turbo_mode 0
gsr set fao_nets VDD
gsr set fao_layers {{metal6 hor <= 29.4}}
gsr set fao_range {{metal6 10 35}}
gsr set noise_reduction 10
\end{lstlisting}
  \item \cmd{mesh optimize -eco <filename>.eco}
  \item 查看结果。
\paragraph{示例 B 结果说明}
收紧 \gsr{noise_reduction} 为正值 10 时，\gsr{fao_range} 须允许更宽导线（本例 metal6 水平线 10--35\,$\mu$m，相对示例 A 的 8--29.4\,$\mu$m 放宽上限）。FAO 在 accurate 模式下搜索全芯片 METAL6 水平导线，将最坏静态 IR 降低约 10\%。优化成功后 metal 用量增加，最坏压降降至约束以内。重跑 \cmd{perform analysis -static} 后对照图~\ref{fig:rh-static-ir-map} 验证整片压降分布。

\paragraph{ECO 导入与回注流程}
\begin{enumerate}
  \item File $>$ Export ECO 写出 \texttt{apache.eco} 或自定义文件名
  \item 在布局布线工具中应用导线/via/pad/decap 变更
  \item 后续 RedHawk 运行：File $>$ Import Design Data 完成 setup 后，File $>$ Import ECO 读入变更
  \item 重跑 \cmd{perform extraction} 与静/动态分析验证收敛
\end{enumerate}

\begin{noteBox}
手动编辑 ECO 文件可能导致电容重复计数或 DRC 违规；应使用 RedHawk 自动生成的 ECO。
\end{noteBox}

\end{enumerate}

\figplaceholder{Figure 7-9}{静态 IR drop 地图}
\label{fig:rh-static-ir-map}

\paragraph{示例 C——局部芯片优化：放宽静态 IR Drop}
条件：设计后期，可接受高于 18\,mV 的压降，但不想修改整片网格；指定 fix window 尺寸，使用 \cmd{mesh fix}。目标：在选定电源网格区域将静态 IR drop 放宽 15\%。
\begin{enumerate}
  \item GSR 设置：
\begin{lstlisting}
gsr set fao_region {2600 2677 6487 5304}
gsr set fao_turbo_mode 0
gsr set fao_nets VDD
gsr set fao_layers { {metal hor <= 29.4} }
gsr set fao_range { {metal6 8 29.4} }
gsr set noise_reduction -15
gsr set fix_window {600 600}
gsr set num_hotspot 1
\end{lstlisting}
  定义角点 (2600,2677) 与 (6487,5304) 的边界区域；FAO 在区域内搜索热点（本例 1 个），以 accurate 模式分析并仅修复 VDD。修改在以热点为中心的 \gsr{fix_window} 内进行，但按金属层约束横跨整片宽/高。搜索宽度 $\leq$29.4\,$\mu$m 的水平 METAL6 线，放宽 15\%；因减少金属用量，宽度范围须小于现有宽度（8--29.4\,$\mu$m）。
  \item \cmd{mesh fix -eco <filename>.eco}
  \item 日志示例：
\begin{lstlisting}
ECO File: fao.static.eco
Region: <2600 2677 6487 5304>
FAO Mode: accurate
Simulation Mode: static
Net: VDD
Noise Reduction: -15%
Mesh Style: uniform
Hot-spots Number: 1
Fix-Window Size: 600x600
Initial Worst Static Noise: 17.578mv
Voltage Noise Constraints: 20.2147mv
Identified 1 hots-spots: (3123.26, 3058.42)
Fixing Window: <2823.26 2758.42 3423.26 3358.42>
Net <VDD>: Layer <metal6>: identified 11 wires
Optimization succeed: no error
Optimized width of layer metal6: 17.9um
Metal Usage Change: metal6 -9.64%
Replace <metal6> mesh grids : net <VDD> : 11 wires replaced.
\end{lstlisting}
  建议修改 11 条 metal6 导线；允许压降升高后金属用量减少。
  \item 重跑静态 IR 分析验证：最终 VDD 静态压降 21\,mV（放宽 15\%），金属资源节省 10\%。
\end{enumerate}

\paragraph{示例 D——局部芯片修复：降低静态 IR Drop}
条件：需要低于 18\,mV 的压降，但仅修改局部区域；使用 \cmd{mesh fix} 与 \gsr{fix_window}。目标：在选定区域将最坏静态 IR drop 降低 10\%。
\begin{enumerate}
  \item 收紧 \gsr{noise_reduction} 时须增大 \gsr{fao_range}：
\begin{lstlisting}
gsr set fao_region {2600 2677 6487 5304}
gsr set fao_turbo_mode 0
gsr set fao_nets VDD
gsr set fao_layers { {metal hor <= 29.4} }
gsr set fao_range { {metal6 10 40} }
gsr set noise_reduction 10
gsr set fix_window {1200 1200}
gsr set num_hotspot 1
\end{lstlisting}
  指定 1200$\times$1200\,$\mu$m 的 fix window；加宽导线范围 10--40\,$\mu$m。
  \item \cmd{mesh fix -eco drop_improv_15.eco}
  \item 日志摘要（图~\ref{fig:rh-static-ir-after-fao}）：
\begin{lstlisting}
Noise Reduction: 10%
Fix-Window Size: 1200x1200
Initial Worst Static Noise(mv): 17.578
Voltage Noise Constraint(mv): 15.8202
Identified 1 hot-spots: (3123.26, 3058.42)
Net <VDD>: Layer <metal6>: identified 17 wires
Optimized width of layer metal6: 37.47um
Metal Usage Change: metal6 +10.46%
Resulting Voltage Noise(mv): 15.8578
\end{lstlisting}
  通过将 17 条 metal6 导线加宽 10.46\% 达到目标压降降低。
  \item 重跑 \cmd{perform analysis -static} 验证。
\end{enumerate}

\figplaceholder{Figure 7-10}{FAO 操作后的静态 IR drop 地图}
\label{fig:rh-static-ir-after-fao}

\paragraph{示例 E——\cmd{mesh optimize -taper} 降低区域静态 IR}
条件：需在特定区域两层上降低压降，使用 \cmd{mesh optimize -taper}。目标：在 \gsr{FAO_REGION} 内对 METAL3、METAL4 做 5\% 静态压降降低，新导线宽度 2--5\,$\mu$m；\opt{-taper} 确保修改仅在 \gsr{fao_region} 内，不横跨整片宽/高。
\begin{enumerate}
  \item GSR 与命令：
\begin{lstlisting}
gsr set fao_nets VDD
gsr set fao_drc 0
gsr set fao_region {x1 y1 x2 y2}
gsr set fao_layers { {METAL4} {METAL3} }
gsr set fao_range { { METAL4 2 5 } { METAL3 2 5 } }
gsr set noise_reduction 5
mesh optimize -taper
export eco mesh_optimize.eco
\end{lstlisting}
  \item FAO 在 \gsr{fao_region} 内搜索可修改的 METAL4/METAL3 导线以达到 5\% 压降降低。
  \item ECO 变更：在区域内删除若干小段导线并以更宽导线替换。
\end{enumerate}

\paragraph{示例 F——\cmd{mesh fix -taper} 降低热点静态 IR}
条件：特定区域两层上热点需更低压降，使用 \cmd{mesh fix -taper}。目标：\gsr{FAO_REGION} 内 5 个热点各降低 2\% 静态压降，METAL3/METAL4 新导线宽度 2--10\,$\mu$m；\opt{-taper} 确保修改仅在各热点 \gsr{fix_window} 内。
\begin{enumerate}
  \item GSR 与命令：
\begin{lstlisting}
gsr set fao_dynamic_mode 1
gsr set fao_nets VDD
gsr set FAO_DRC 0
gsr set fao_region { x1 y1 x2 y2 }
gsr set fao_layers { {METAL4} {METAL3} }
gsr set fao_range { { METAL4 2 10} { METAL3 2 10 } }
gsr set noise_reduction 2
gsr set fix_window { 600 600 }
gsr set num_hotspot 5
mesh fix -taper
export eco mesh_fix.eco
\end{lstlisting}
  \item FAO 在区域内识别 5 个热点，在各 600$\times$600\,$\mu$m fix window 内寻找至少 2\% 压降降低的网格修改。
  \item ECO 变更：在 5 个热点周围 fix window 内删除小段导线并以更宽导线替换。
\end{enumerate}

\section{DvD 与时序的自动修复与优化（FAO）}
\subsection*{Automated Fixing and Optimization (FAO) for DvD and Timing}

\subsection{概述}
\subsubsection*{Overview}

RedHawk 修复与优化动态压降与时序问题的流程见图~\ref{fig:rh-fao-flow}。三种 FAO 运行准则：
\begin{itemize}
  \item 最坏动态压降（DvD）的 instance（热点 instance）
  \item DvD 导致最高 delta slack 与 delta delay 的全部 instance
  \item 关键路径上 DvD 导致最高 delta slack 与 delta delay 的 instance
\end{itemize}

\figplaceholder{Figure 7-11}{FAO 流程}
\label{fig:rh-fao-flow}

根据 GSR/FAO 关键字 \gsr{FAO_OBJ}，程序创建有序 instance 列表：纯 DvD 运行按最坏 DvD 排序；时序相关运行按 DvD 影响严重程度排序。

\gsr{FAO_OBJ} 语法：
\begin{lstlisting}
FAO_OBJ [ DVD | TIMING ?<slack threshold>? |
      PATH ?<slack threshold>? ]
\end{lstlisting}
\begin{itemize}
  \item \texttt{DVD}：选择压降最高的 instance
  \item \texttt{TIMING <slack\_threshold>}：选择 DvD 导致高 delta delay/slack 的 instance（默认 slack 阈值 0）。使用前须 \cmd{import sdf xxx.sdf}，否则 swap 候选为零。\gsr{fao_region} 仍用于过滤。
  \item \texttt{PATH <slack\_worse\_than>}：关键时序路径上 DvD 影响显著的 instance
\end{itemize}

可用 DvD 值、instance 总数或感兴趣区域进一步过滤目标 instance。

\subsection{Decap FAO 准备}
\subsubsection*{Preparation for Decap FAO}

\begin{enumerate}
  \item 在 GSR 中用 \gsr{DECAP_CELL} 定义 decap 列表：
\begin{lstlisting}
DECAP_CELL
{
<decap_cellname> ? <width_um> <height_um> <C_pF> <R_ohm>
<metal_layer_name> ?
...
}
\end{lstlisting}
  \item 运行 APL 表征，生成含 R、C、漏电的 \texttt{<cell>.cdev}。
  \item 确保 LEF 中物理信息与 P/G pin 定义完整。
\end{enumerate}
\paragraph{DECAP\_CELL 与 LEF 物理定义示例}
\gsr{DECAP_CELL} 通常仅需列出 decap 单元名，其余信息在 LEF 中定义。LEF 中 decap 宏须含完整 P/G pin 定义，例如：
\begin{lstlisting}
MACRO cell_name
    CLASS <class_type> ;
    ORIGIN <value> ;
    SIZE <value> BY <value> ;
    PIN GND
        USE GROUND ;
        PORT
            LAYER <layer_name> ;
            RECT <x1> <y1> <x2> <y2> ;
        END
    END GND
    PIN VDD
        USE POWER ;
        PORT
            LAYER <layer_name> ;
            RECT <x1> <y1> <x2> <y2> ;
        END
    END VDD
END cell_name
\end{lstlisting}


\subsection{Decap 修改操作}
\subsubsection*{Decap Modification Operations}

RedHawk FAO 可调查使用去耦电容修复动态压降与时序。

\paragraph{去耦电容修改命令}
\begin{itemize}
  \item \cmd{decap advise}——快速评估可放置 decap 总量（\opt{-place} 生成 dpf 文件）
  \item \cmd{decap fill}——多种方式放置 decap：targeted、hot instance area、region、prefill、uniform
  \item \cmd{decap remove}——移除指定区域内无效 decap（默认峰值电流 $\leq$30\,$\mu$A）
\end{itemize}

\paragraph{Decap 修改约束与外部程序接口}
\begin{itemize}
  \item \cmd{decap advise -place}、\cmd{decap fill}、\cmd{cell swap} 遵守 DEF \texttt{+FIXED}
  \item \gsr{FAO_DRC_PL_OBS} 为 1 时遵守 DEF placement blockage
  \item decap 含 metal1 结构时，\gsr{FAO_DRC_OBS} 为 1 时检测与 metal1 信号布线 DRC
  \item GSR 所列 decap 在 \cmd{cell swap} 中若指定 \opt{-fixed_inst/cell} 则不移动
  \item tie high/low 单元默认不移动（可用 \opt{-include_thtl}；布线前可移动）
\end{itemize}

\paragraph{控制 Decap 修改的 GSR 关键字}
\begin{enumerate}
  \item \gsr{NOISE_REDUCTION}——目标压降降低（必需）
  \item \gsr{NOISE_LIMIT}——压降改善百分比
  \item \gsr{FAO_REGION}——修复区域
  \item \gsr{FIX_WINDOW}——热点周围区域（默认 40$\times$40\,$\mu$m）
  \item \gsr{NUM_HOTINST}——修复的最高压降区域数（默认 10）
  \item \gsr{CAP_LIMIT}——新增 decap 总电容上限
  \item \gsr{LEAKAGE_LIMIT}——新增 decap 总漏电上限
  \item \gsr{FAO_DECAP_OVERLAP}——仅 \cmd{decap advise -place}
  \item \gsr{DECAP_DENSITY}——每种合法 decap 的放置倍数
\end{enumerate}

\subsection{高压降区域的附加修复工具}
\subsubsection*{Additional Tools for Fixing High Voltage Drop Areas}

\paragraph{使用 \cmd{route fix} 的补充电源布线}
高压降节点可通过 \cmd{route fix} 布线到邻近低压降节点，形成并联路径均衡压降：
\begin{lstlisting}
route fix [-from { x1 y1 x2 y2 <layer>} | -from_pt { xf yf <layer>} ]
[-to {x1 y1 x2 y2 <layer>} | -to_pt { xt yt <layer>} ] -width <float>
\end{lstlisting}
查找最短 DRC 合法路由；\cmd{export eco} 定义新增导线/via。新 via 可能记录在 \texttt{adsRpt/faoVIAS.def}。用于从较强点（压降较小）改善到较弱点的电源布线（据 PG-Aware 报告识别）。

\paragraph{热点 Instance 的移动或单元替换}
可通过将「热点」instance 移至同网格上压降更好的节点，或与其他 instance 交换以均衡电流需求。

\cmd{cell swap} 流程概要：
\begin{enumerate}
  \item 设置 \gsr{FAO_OBJ}：\texttt{DVD}、\texttt{TIMING} 或 \texttt{PATH}
  \item 选择 \gsr{fix_window} 大小（默认 40$\times$40\,$\mu$m）
  \item 选择 DvD 准则：\opt{-eff_vdd_tw}（默认）、\opt{-max_vdd_tw}、\opt{-min_vdd_tw}、\opt{-min_vdd_all}
  \item 用 \opt{-fixed_cells}、\opt{-fixed_insts} 等排除不可交换单元
  \item 调用 \cmd{cell swap}
  \item FAO 生成候选列表，尝试移动与交换
  \item 自动准备 ECO 报告；可用 \opt{-report}、\opt{-plot} 查看结果
\end{enumerate}

\cmd{cell swap} 详细语法：
\begin{lstlisting}
cell swap ?-eco <file_name>?
      ?-decap_cells {cell1 cell2 ...}?
      ?-fixed_cells {cell1 cell2 ...}?
      ?-fixed_insts {inst1 inst2 ...}?
      ?-removable_cells {cell1 cell2 ...}?
      ?-glob|-regexp|-exact?
      ?-include_clock ?
      ?-inst_list {inst1 inst2 ...}?
      ?-inst_file <file_name>?
      ?-search_only?
      ?-eff_vdd_tw? ?-max_vdd_tw?
      ?-min_vdd_tw? ?-min_vdd_all?
      ?-report? ?-plot?
\end{lstlisting}

相关 GSR 关键字：\gsr{fao_region}、\gsr{fao_nets}、\gsr{noise_limit}、\gsr{num_hotinst}、\gsr{fix_window}、\gsr{fao_verbose}。

\figplaceholder{Figure 7-12}{Cell Swapping 报告示例}
\paragraph{\cmd{cell swap} 选项详述}
\begin{itemize}
  \item \opt{-eco <file\_name>}：写出描述设计变更的 ECO 文件名
  \item \opt{-decap\_cells \{cell1 ...\}}：可与热点 instance 交换的 decap master 单元；默认任意无逻辑连接单元可交换。\opt{-glob} 允许通配符；\opt{-regexp} 允许正则；\opt{-exact} 精确单元名
  \item \opt{-fixed\_cells \{...\}}：禁止参与交换的 master 单元
  \item \opt{-fixed\_insts \{...\}}：禁止参与交换的 instance
  \item \opt{-removable\_cells \{...\}}：必要时可从设计中移除的 decap 单元（默认不可移除）；同样支持 \opt{-glob/-regexp/-exact}
  \item \opt{-include\_clock}：允许交换时钟单元
  \item \opt{-inst\_list \{...\}} / \opt{-inst\_file <file>}：仅对列出的 instance 执行交换
  \item \opt{-search\_only}：仅生成候选 instance 列表（默认 \texttt{cell\_swap\_<time>.list}），不执行交换；可用 \opt{-inst\_file} 改名
  \item DvD 准则：\opt{-eff\_vdd\_tw}（TW 上平均有效电压，默认）、\opt{-max\_vdd\_tw}、\opt{-min\_vdd\_tw}、\opt{-min\_vdd\_all}（整周期最小有效电压）
  \item \opt{-report}：生成 \texttt{adsRpt/cell\_swap.rpt} 文本表；\opt{-plot}：生成 \texttt{adsRpt/cell\_swap.hist} 直方图
\end{itemize}

\paragraph{Cell Swapping 报告日志示例}
\begin{lstlisting}
*******************************************************************************
**** FAO Cell Swapping Report
*******************************************************************************
Cell Swapping Result Summary
   Total swapping instance number : 2182
   Actually Moved instance number : 436
   Actually Deleted instance number: 0
   maximum voltage change (mv)                      : 171
   minimum voltage change (mv)                      : -18
**** Please see adsRpt/cell_swap.rpt for detail DvD change.
**** To view histogram of DvD change, please execute xgraph
adsRpt/cell_swap.hist.
MEMORY USAGE: 2265 MBytes
TOTAL CPU TIME: 1 hrs 17 mins 52 secs
WALLTIME: 1 hrs 27 mins 56 secs.
*******************************************************************************
\end{lstlisting}

\paragraph{控制 \cmd{cell swap} 的 GSR 关键字}
\begin{itemize}
  \item \gsr{fao\_region}：含高功耗 instance 的目标分析区域
  \item \gsr{fao\_nets}：多 VDD 设计时指定热点 instance 关联的 VDD 网
  \item \gsr{noise\_limit}：对最差热点列表，低于该压降百分比阈值的热点不参与交换（例：5\% 表示压降 $<$5\% 的热点忽略）
  \item \gsr{num\_hotinst}：区域内交换的最差热点 instance 数
  \item \gsr{fix\_window}：以各热点为中心的可移动/交换区域尺寸（默认 40$\times$40\,$\mu$m）
  \item \gsr{fao\_verbose}（0/1）：控制输出消息量（默认 0）
\end{itemize}


\subsection{DvD 的 FAO 示例}
\subsubsection*{Examples of FAO for DvD}

\paragraph{示例 G——全芯片 DvD 降低：非重叠 Decap 与网格修复}
初始最坏 DvD 306\,mV，目标降低 10\%（至 276\,mV）：
\begin{enumerate}
  \item \cmd{gsr set noise_reduction 10}
  \item \cmd{decap advise -place}
  \item 查看目标 DvD 报告（TW 上最差平均 DvD 百分比）
  \item 重跑提取与动态分析验证；\cmd{decap report}、\cmd{export eco decap.1.eco}
  \item 若未达目标，对剩余改善量设置 \gsr{noise_reduction} 并运行 \cmd{mesh optimize} 或 \cmd{mesh fix}
\end{enumerate}

\paragraph{示例 H——全芯片 DvD 降低：Decap 重叠}
设置 \gsr{fao_decap_overlap 1} 允许在无合法空间处也插入 decap，示例总体 DvD 改善约 7.73\%。
\paragraph{示例 G 详细步骤与日志}
\begin{enumerate}
  \setcounter{enumi}{2}
  \item FAO 评估期间报告目标 DvD，例如：
\begin{lstlisting}
|Target Worst Average DvD percent over TW: 24.1895%(0.27576)(v.s. Vdd=1.14)
\end{lstlisting}
  表示目标 276\,mV 对应 Vdd 的 24.2\%；测量值为 TW 上最差平均 DvD（STA 给出的门可能开关时间窗）。
  \item FAO 自动迭代放置 decap 并报告，例如首轮：
\begin{lstlisting}
Placed decap cells for 29 hot instances.
Added decap 5.092 pF
\end{lstlisting}
  第二轮累计：
\begin{lstlisting}
Placed decap cells for 29 hot instances.
Added decap 204.718 pF
\end{lstlisting}
  \item 重跑提取与动态分析验证：
\begin{lstlisting}
perform extraction -power -ground -c
setup pad -power -r 0.005 -ground -r 0.005
setup wirebond -power -r 0.001 -l 1000 -ground -r 0.001 -l 1000
setup package -r 0.001 -l 250 -c 5
perform analysis -vectorless
\end{lstlisting}
  本例 DvD 改善至 284\,mV（目标 276\,mV）。
  \item \cmd{decap report} 完整报告示例：
\begin{lstlisting}
*************************************************************
****   FAO Decap Report on top Hot Instances
*************************************************************
Inserted Decap Inst Num : 3579
Inserted Decap Total      : 209.81 pF
V0/V1 : Average (VDD - GND) over TW before/after decap placement.
DvD0/DvD1 : Dynamic Voltage Drop (VDD - V0/V1) before/after decap
placement. Improvement % : (DvD0 - DvD1)/DvD0 %.
inst    V0    V1     DvD0     DvD1    Improvement%
-------------------------------------------------
INFO(FAO-75): inst194510 : 0.8343 0.8568 0.3057 0.2832 7.36017%
INFO(FAO-75): inst117100 : 0.8346 0.8574 0.3054 0.2826 7.46563%
INFO(FAO-75): inst76925 : 0.8395 0.8607 0.3005 0.2793 7.05491%
--------------------------------------
New Top Hot Instances:
inst    &   avgVoltage_TW
--------------------------------------
inst194510 0.8568
inst117100 0.8574
inst76925 0.8607
--------------------------------------
Overall DvD improvement = 7.3601 % v.s. initial DvD
\end{lstlisting}
  \item 在 ECO 中高亮新增 decap：\cmd{select add "decap*" -glob -linewidth 3}
  \item 将剩余改善量（本例 2.64\%）赋给网格修复：\cmd{gsr set noise\_reduction 2.64}，再运行 \cmd{mesh optimize} 或 \cmd{mesh fix}
\end{enumerate}

\paragraph{示例 H 完整 decap 重叠报告}
\begin{lstlisting}
**********************************************************
****      FAO Decap Report on top Hot Instances
**********************************************************
Inserted Decap Inst Num : 305
Inserted Decap Total             : 229.938 pf
V0/V1 : Average (VDD - GND) over TW before/after decap fix.
DvD0/DvD1:Dynamic Voltage Drop (VDD - V0/V1) before/after decap fix.
Improvement % : (DvD0 - DvD1)/DvD0 %.
     inst                       V0     V1 DvD0 DvD1 Improvement%
--------------------------------------------------------
INFO(FAO-75): inst117100 : 0.8346 0.8573 0.3054 0.2827 7.43288%
INFO(FAO-75): inst194510 : 0.8344 0.8609 0.3056 0.2791 8.67146%
INFO(FAO-75): inst76925 : 0.8396 0.8623 0.3004 0.2777 7.55658%
-------------------------------------------------------
New Top Hot Instances:
inst       &    avgVoltage_TW
-------------------------------
inst117100 0.8573
inst194510 0.8609
inst76925 0.8623
-------------------------------
Overall DvD improvement = 7.73499 % v.s. initial DvD
\end{lstlisting}


\section{使用 ECO 命令保存设计变更}
\subsection*{Saving Design Changes with the ECO Command}

达到压降目标后，可用 GUI File $>$ Import ECO 读入 mesh 与 decap 操作生成的 ECO 文件（ASCII 文本，便于导出到布线工具）。

\subsection{写出 ECO 文件}
\subsubsection*{Writing an ECO File}

File $>$ Export ECO 将近期设计变更导出到工作目录。弹出表单查询文件名，默认 \texttt{apache.eco}。TCL 见附录 D \cmd{export eco}。

\subsection{读取 ECO 文件}
\subsubsection*{Reading an ECO File}

File $>$ Import Design Data 读入设计并完成 database setup 后，可用 File $>$ Import ECO。ECO 或 \texttt{apache.eco} 中的 what-if 变更将纳入 IR drop 与 EM 分析。

\subsection{ECO 文件格式定义}
\subsubsection*{ECO File Format Definition}

\begin{noteBox}
以下语法供理解 RedHawk 自动生成的 ECO 内容。\textbf{请勿}手动创建或编辑 ECO 文件。
\end{noteBox}

示例表格式 \texttt{*.eco}：
\begin{lstlisting}
Keyword  Obj  Obj name  P/G  Layer/via  Rotation  Coordinates  Direction
add      pad  pvdd1     Vdd  metal1     None      x_new y_new    None
delete   pad  pvdd2     Gnd  metal1     None      x_old y_old    None
add      via  via67_new Vdd1 via7_1     <code>    x_new y_new    None
add      wire wire1     Vdd2 metal7     None      x1 y1 x2 y2    horizontal|vertical
add      wire45 wireCD  Vdd  metal7     None      x3 y3 ... x6 y6  None
\end{lstlisting}

各列含义：
\begin{itemize}
  \item Keyword：\texttt{add} 或 \texttt{delete}
  \item Obj：\texttt{pad}、\texttt{via}、\texttt{wire}、\texttt{wire45}
  \item Obj name：对象唯一名
  \item P/G Name：pad 为 VDD/GND；导线/via 为网络名
  \item layer/via：LEF 中金属层或 via 名
  \item Coordinates：pad 中心；via 旋转码（-1 无旋转，0--6 RedHawk 旋转码）；导线 bounding box；wire45 为梯形四角（逆时针）
  \item Direction：\texttt{add wire} 时为 vertical 或 horizontal
\end{itemize}

\begin{noteBox}
ECO 添加的导线不会被 RedHawk 合并；同层同网重叠会导致电容重复计数、电阻偏小，应尽量减少重叠。
\end{noteBox}

\subsection{供布局布线工具使用的 ECO 转换}
\subsubsection*{ECO File Translation for Use by Place and Route Tools}

ECO 格式可后处理为 EDA 布局布线工具可读格式，确保 RedHawk 与 P\&R 工具间数据流顺畅。

\section{修复与优化命令参考}
\subsection*{Fixing and Optimization Command Reference}

\subsection{Mesh 优化命令说明}
\subsubsection*{Descriptions of Mesh Optimization Commands}

Mesh 与 ring 命令（按字母序）：
\begin{lstlisting}
mesh [ add| delete| fix| generate| optimize| snscalc |
    sub_grid | set_width | vias ] ? args ?
ring add ? args ?
ring delete <ring_name>
\end{lstlisting}

表~\ref{tab:mesh-commands} 汇总网格优化与修复命令（对应原书表 7-1）。各命令关联的 GSR 关键字见附录 C「Grid Fixing and Optimization Keywords」（C-724）。

\begin{longtable}{@{}p{2.4cm}p{10.6cm}@{}}
\caption{网格优化与修复命令（表 7-1）}
\label{tab:mesh-commands}\\
\toprule
命令 & 说明与主要选项 \\
\midrule
\endfirsthead
\multicolumn{2}{l}{\textit{（续表~\ref{tab:mesh-commands}）}}\\
\toprule
命令 & 说明与主要选项 \\
\midrule
\endhead
\bottomrule
\endfoot
\cmd{mesh add} &
在电源网格上添加指定 mesh。关联 GSR：无。\\
& \opt{-vertical|-horizontal}：方向；\opt{-window \{x1 y1 x2 y2\}}：矩形区域；\\
& \opt{-layer <layer>}：层名；\opt{-offset <um>}：自窗口左/底边偏移；\\
& \opt{-space <um>}：\opt{-net} 中各电源网间距；\opt{-pitch <um>}：pitch；\\
& \opt{-width <um>}：线宽；\opt{-novia} 或 \opt{-viatop/-viabottom}：叠层 via；\\
& \opt{-net \{net1 ...\}}：网络；\opt{-clip_cell \{macro...\}}：宏内禁布金属防短路；\\
& \opt{-exclude_regions \{\{x1 y1 x2 y2\}...\}}：排除区域。\\
\cmd{mesh snscalc} &
计算所选层/网/区域内导线宽度变化的相对压降灵敏度（mv/$\mu$m）。\\
& 关联 GSR：\gsr{fao_region}、\gsr{fao_nets}、\gsr{fao_layers}。\\
& \opt{-taper}：仅在 \gsr{fao_region} 内调整，非整片宽/高。\\
\cmd{mesh delete} &
删除指定网格段。关联 GSR：无。\\
& \opt{-vertical|-horizontal}、\opt{-window}、\opt{-layer}、\opt{-net}；\\
& \opt{-width \{?<Oper>? <w-1> ...\}}、\opt{-length \{?<Oper>? ...\}}（Oper 为 =,>,<,$\geq$,$\leq$）；\\
& \opt{-out_of_ratio \{n m\}}：每 m 条线删 n 条；\opt{-all}；\opt{-exclude_regions}。\\
\cmd{mesh fix} &
基于热点的网格宽度修复。关联 GSR：\gsr{fao_region}、\gsr{mesh_search}、\gsr{fao_layers}、\gsr{fao_turbo_mode}、\gsr{fao_dynamic_mode}、\gsr{fao_nets}、\gsr{fao_range}、\gsr{fao_width_cnstr}、\gsr{noise_reduction}、\gsr{fix_window}、\gsr{num_hotspot}。\\
& \opt{-taper}：仅在 \gsr{FAO_REGION} 或 \gsr{fix_window} 内调整；\opt{-eco <file>}：写出 ECO。\\
\cmd{mesh optimize} &
均匀网格修改以调查/修复压降。关联 GSR 同 \cmd{mesh fix}（不含 \gsr{fix_window}）。\\
& \opt{-taper}、\opt{-eco <file>}。\\
\cmd{mesh sub_grid} &
在 \gsr{fao_region} 内按 \gsr{fao_sub_grid_nets}、\gsr{fao_sub_grid_spec} 添加子网格，目标压降 \gsr{noise_reduction}，宽度范围 \gsr{fao_range}。\\
& \opt{-eco <file>}。\\
\cmd{mesh vias} &
添加最优尺寸叠层 via 或报告/删除缺失 via。关联 GSR：\gsr{FAO_LAYERS}、\gsr{FAO_NETS}。\\
& 添加：\opt{-toplayer/-bottomlayer}、\opt{-viarule}、\opt{-window|-inst}、\opt{-net}、\\
& \opt{-topwidth/-bottomwidth}、\opt{-viamodel}、\opt{-replace}、\opt{-between_layer}、\opt{-eeq_net}、\opt{-pt_file}；\\
& \opt{-vr_percent_x/y}、\opt{-vr_x/-vr_y}：via 阵列尺寸约束。\\
& 报告缺失：\opt{-report_missing}、\opt{-exclude_stack_via}、\opt{-check_valid_vias}、\\
& \opt{-exclude_inst_region/-exclude_cell_region}、\opt{-exclude_cell/-exclude_instance}、\\
& \opt{-inst/-cell/-window}、\opt{-exclude_regions}、\opt{-sort_by_hot_inst}、\opt{-reverse_sort_v}、\\
& \opt{-gds}、\opt{-threshold}（默认 1\,mV，\texttt{-1} 禁用）、\opt{-min_width}、\opt{-fao_layers}、\\
& \opt{-marker}、\opt{-no_partial_overlap}、\opt{-pitch}、\opt{-from_edge}、\opt{-x/y_pitch}、\opt{-x/y_offset}、\opt{-adjust_offset}、\opt{-ignore_inter_met}。\\
& 删除：\opt{-delete -window ... -viamodel ... -cut_layer}。\\
\cmd{mesh set_width} &
修改指定层 mesh 线宽。关联 GSR：\gsr{fao_region}、\gsr{fao_sub_grid_nets}、\gsr{fao_layers}。\\
& 语法：\texttt{\{<layer> <new_width> ?<expand_dir>? \}}，\texttt{expand_dir} 为 left/right/top/bot。\\
\cmd{mesh generate} &
按设计约束文件（.dcf）生成原型网格。\\
\cmd{ring add} &
添加电源/地 ring。\opt{-name}、\opt{-window}；\\
& \texttt{\{-top|bottom|left|right -layer -offset -width -space -net \{...\}\}}。\\
\cmd{ring delete} &
删除指定 ring：\texttt{<ring_name>}。\\
\end{longtable}

\subsection{Decap 修改命令说明}
\subsubsection*{Descriptions of Decap Modification Commands}

\begin{lstlisting}
decap [ advise | place | fill | remove | report ] ? arg ?
cell swap ? arg?
\end{lstlisting}

表~\ref{tab:decap-commands} 汇总 decap 与 cell swap 命令（对应原书表 7-2）。

\begin{longtable}{@{}p{2.4cm}p{10.6cm}@{}}
\caption{Decap 修改命令（表 7-2）}
\label{tab:decap-commands}\\
\toprule
命令 & 说明与主要选项 \\
\midrule
\endfirsthead
\multicolumn{2}{l}{\textit{（续表~\ref{tab:decap-commands}）}}\\
\toprule
命令 & 说明与主要选项 \\
\midrule
\endhead
\bottomrule
\endfoot
\cmd{decap report} &
生成 decap 报告（\texttt{adsRpt/faoOverlapDecap}：总数、电容、DvD 改善）。\\
& \opt{-overlap}：重叠 decap 统计；DvD 准则：\opt{-eff_vdd_tw}（默认）、\opt{-max_vdd_tw}、\opt{-min_vdd_tw}、\opt{-min_vdd_all}。\\
\cmd{decap drc} &
DRC 报告；\opt{-fix} 修复用户放置 decap 的 pin 与信号/屏蔽布线 DRC；\opt{-o <file>}。\\
\cmd{decap advise} &
评估可放置 decap 量；\opt{-place} 生成 \texttt{*.dpf}；\opt{-dpf}、\opt{-distribute}、DvD 准则选项、\opt{-inst_file/-inst}、\opt{-repeat [-iter n]}、\opt{-mmx}。\\
& 关联 GSR：\gsr{fao_region}、\gsr{noise_reduction}、\gsr{noise_limit}、\gsr{num_hotinst}、\gsr{cap_limit}、\gsr{fao_decap_overlap}、\gsr{leakage_limit}、\gsr{decap_density}。\\
\cmd{decap place} &
按 advise 结果放置 decap。\opt{-dpf}（默认 \texttt{adsRpt/fao.dpf}）、\opt{-eco}。\\
\cmd{decap qor} &
评估 FAO 插入 decap 的 QoR。\opt{-i} 输入实例列表；\opt{-o} 输出目录。关联 \gsr{decap_fill}。\\
\cmd{decap fill} &
在可用空间填充 decap。\opt{-repeat}、\opt{-density_step}、\opt{-eco_name_pattern}、\opt{-no_adjust_fix_window}、\opt{-clone_cells}、\opt{-prefill/-prefill_only/-no_ff_move/-ignore_fixed}、\opt{-uniform <pct>}、\opt{-fao_region/-hot_area}、DvD 准则、\opt{-inst_file/-inst}、\opt{-fixed_insts/-fixed_cells}、\opt{-replace_filler}。\\
\cmd{decap remove} &
移除无效 decap（默认峰值电流 $\leq$30\,$\mu$A）。\opt{-imax_less_than}、\opt{-dvmax_less_than}、\opt{-fao_region/-hot_area}、\opt{-user_decap}、DvD 准则、\opt{-ignore_cap_limit/-ignore_leakage_limit}、\opt{-no_decap_report}、\opt{-clone_cells}、\opt{-all}。\\
\cmd{cell swap} &
移动/交换热点 instance 与 decap 或低功耗单元。\opt{-eco}、\opt{-decap_cells}（\opt{-glob/-regexp/-exact}）、\opt{-fixed_cells/-fixed_insts/-removable_cells}、\opt{-include_clock}、\opt{-inst_file/-inst}、\opt{-search_only}、DvD 准则、\opt{-report/-plot}、\opt{-swappable_cells}、\opt{-ff_only}、\opt{-pre_proc}、\opt{-to_scan_in_only}、\opt{-include_thtl}、\opt{-drc_m2}。\\
\cmd{route fix} &
点对点补充电源布线（表~\ref{tab:route-fix}）。\\
\cmd{fao add pad} &
批量添加 pad（见「在指定区域添加一组电源/地 Pad」）。\\
\cmd{fao undo/redo} &
撤销/重做 FAO 变更。\\
\cmd{export eco} &
导出 ECO 文件。\\
\end{longtable}

\subsection{补充压降修复命令}
\subsubsection*{Supplementary Voltage Drop Fixing Commands}

\begin{table}[htbp]
\centering
\caption{Route Fix 命令（表 7-3）}
\label{tab:route-fix}
\small
\begin{tabular}{@{}p{2.5cm}p{10cm}@{}}
\toprule
命令 & 说明 \\
\midrule
\cmd{route fix} &
查找从指定起点到终点的最短 DRC 合法电源/地布线与 via。\\
& \opt{-from \{x1 y1 x2 y2 <layer>\}|-from_pt \{xf yf <layer>\}}：起点区域或点；\\
& \opt{-to \{...\}|-to_pt \{xt yt <layer>\}}：终点；\opt{-width <um>}：线宽。\\
& \cmd{export eco} 定义新增导线/via；新 via 可能记录在 \texttt{adsRpt/faoVIAS.def}。\\
\bottomrule
\end{tabular}
\end{table}

\subsection{Mesh 命令选项补充（表 7-1 扩展）}
\subsubsection*{Extended Mesh Command Options (Table 7-1)}

\paragraph{\cmd{mesh add} 完整选项}
\begin{itemize}
  \item \opt{-vertical|-horizontal}：添加垂直或水平 mesh
  \item \opt{-window \{x1 y1 x2 y2\}}：矩形区域
  \item \opt{-layer <name>}：金属层
  \item \opt{-offset <um>}：自窗口左/底边偏移
  \item \opt{-space <um>}：\opt{-net} 中各电源网间距
  \item \opt{-pitch <um>}：pitch
  \item \opt{-width <um>}：线宽
  \item \opt{-novia} 或 \opt{-viatop/-viabottom}：无 via 或叠层至指定层
  \item \opt{-net \{net1 ...\}}：关联网
  \item \opt{-clip_cell \{macro...\}}：宏边界内禁布金属防短路
  \item \opt{-exclude_regions \{\{x1 y1 x2 y2\}...\}}：排除区域
\end{itemize}

\paragraph{\cmd{mesh delete} 宽度/长度过滤}
\opt{-width \{?<Oper>? <w> ...\}}、\opt{-length \{?<Oper>? ...\}}，Oper 为 \texttt{=}, \texttt{>}, \texttt{<}, \texttt{$\geq$}, \texttt{$\leq$}。\opt{-out_of_ratio \{n m\}}：每 m 条线删 n 条。\opt{-all}：删除所有匹配段。

\paragraph{\cmd{mesh vias} 报告缺失 via}
\opt{-report_missing} 配合 \opt{-threshold}（默认 1\,mV，\texttt{-1} 禁用）、\opt{-sort_by_hot_inst}、\opt{-gds} 输出 GDS marker、\opt{-pitch}、\opt{-x/y_pitch}、\opt{-x/y_offset}、\opt{-adjust_offset}、\opt{-ignore_inter_met} 等控制阵列尺寸与排序。

\paragraph{\cmd{ring add}}
\opt{-name}、\opt{-window}；\texttt{\{-top|bottom|left|right -layer -offset -width -space -net \{...\}\}} 定义 ring 各边。

\paragraph{FAO 流程子步骤（时序感知）}
时序相关 FAO 在 DvD 热点基础上按 delta slack/delay 排序，可选子网格（sub-gridding）、重定线宽、via 插入、cell swapping、decap 插入等组合策略（图~\ref{fig:rh-fao-flow}）。

\subsection{支持 FAO 的 GSR 关键字}
\subsubsection*{GSR Keywords Supporting FAO Functionality}

\gsr{FAO_OBJ}、\gsr{FAO_NETS}、\gsr{FAO_REGION}、\gsr{FAO_LAYERS}、\gsr{FAO_RANGE}、\gsr{NOISE_REDUCTION}、\gsr{NOISE_LIMIT}、\gsr{FIX_WINDOW}、\gsr{NUM_HOTINST}、\gsr{NUM_HOTSPOT}、\gsr{CAP_LIMIT}、\gsr{LEAKAGE_LIMIT}、\gsr{DECAP_CELL}、\gsr{DECAP_DENSITY}、\gsr{FAO_DECAP_OVERLAP}、\gsr{FAO_DRC}、\gsr{FAO_TURBO_MODE}、\gsr{FAO_DYNAMIC_MODE} 等修改 \cmd{mesh optimize}、\cmd{mesh fix}、\cmd{route fix}、\cmd{decap advise}/\cmd{place}、\cmd{cell swap} 的行为。详见附录 C「FAO General Keywords」「Grid Fixing and Optimization Keywords」「Decap Optimization Keywords」（C-722、C-724、C-728）。

\subsection{命令与 GSR 关键字语法约定}
\subsubsection*{Command and GSR Keyword Syntax Conventions}

\begin{itemize}
  \item \texttt{<x>}：须指定的变量名或值
  \item \texttt{[ a | b | c ]}：须选 a、b、c 之一
  \item \texttt{? x ?}：可选参数
  \item \texttt{\{ a b c \}}：a、b、c 均须指定
  \item \texttt{\{ x \} ...}：可重复添加 \texttt{\{x\}} 元素
  \item \texttt{+ - * /}：加减乘除算术运算符
\end{itemize}

\begin{seeAlsoBox}
\subsection{FAO GSR 关键字详述}
\subsubsection*{Detailed FAO GSR Keywords}

\paragraph{\gsr{FAO\_DYNAMIC\_MODE}}
选择静态 IR 或动态 DvD 分析作为 FAO 优化依据。\texttt{[0|1]}，默认 0（静态）。

\paragraph{\gsr{FAO\_TURBO\_MODE}}
选择 accurate 或 turbo 搜索模式。默认 on（turbo）；设为 0 使用 accurate 模式，精度更高但更慢。

\paragraph{\gsr{FAO\_RANGE}}
指定各金属层允许的新导线宽度范围（$\mu$m）：
\begin{lstlisting}
FAO_RANGE { {metal6 8 29.4} {metal5 4 12} }
\end{lstlisting}
放宽压降（负 \gsr{noise\_reduction}）时范围须小于现有宽度；收紧压降时须增大上限。

\paragraph{\gsr{FAO\_LAYERS}}
指定参与优化的层与方向约束：
\begin{lstlisting}
FAO_LAYERS { {metal6 hor <= 29.4} {metal5 ver} }
\end{lstlisting}
支持 \texttt{hor/ver}、宽度比较运算符（\texttt{<=}、\texttt{>=} 等）。

\paragraph{\gsr{MESH\_SEARCH}}
指定各层压降降低百分比分配（默认各选定层均分）：
\begin{lstlisting}
MESH_SEARCH { {metal6 60} {metal5 40} }
\end{lstlisting}

\paragraph{\gsr{FAO\_WIDTH\_CNSTR}}
定义层间宽度关系约束（可选）。

\paragraph{\gsr{FAO\_SUB\_GRID\_NETS} / \gsr{FAO\_SUB\_GRID\_SPEC}}
\cmd{mesh sub\_grid} 所用：子网格电源网列表及物理特性（pitch、方向等）。

\paragraph{\gsr{FAO\_DRC} / \gsr{FAO\_DRC\_OBS} / \gsr{FAO\_DRC\_PL\_OBS}}
控制 FAO 是否执行 DRC 检查：\gsr{FAO\_DRC} 总开关；\gsr{FAO\_DRC\_OBS} 为 1 时 decap 与 metal1 信号布线做 DRC；\gsr{FAO\_DRC\_PL\_OBS} 为 1 时遵守 DEF placement blockage。

\paragraph{\gsr{DECAP\_FILL}}
与 \cmd{decap qor} 配合，指定 decap 填充策略关键字（见附录 C）。

各命令完整选项与 GSR 语法见附录 C「GSR File Keywords」与「Grid Fixing and Optimization Keywords」。
\end{seeAlsoBox}


\subsection{Decap 命令选项扩展说明}
\subsubsection*{Extended Decap Command Reference}

\paragraph{\cmd{decap advise} 完整流程}
\begin{enumerate}
  \item 设置 \gsr{noise_reduction}、\gsr{fao_region}、\gsr{num_hotinst} 等
  \item \cmd{decap advise} 快速估算可放置总量
  \item \cmd{decap advise -place} 生成 \texttt{*.dpf} 放置文件
  \item \cmd{decap place -dpf <file>} 或 \cmd{decap fill} 执行放置
  \item \cmd{decap report} 查看改善；\cmd{export eco} 导出变更
\end{enumerate}

\paragraph{\cmd{decap fill} 模式详解}
\begin{itemize}
  \item \textbf{targeted}：每个热点 instance 周围矩形区域内放置
  \item \textbf{hot instance area}： encompass 一组热点的矩形区域
  \item \textbf{region}：用户指定 \gsr{fao_region}
  \item \textbf{prefill}：在热点旁腾出行空间后相邻放置
  \item \textbf{uniform}：\opt{-uniform <pct>} 在所有标准单元行均匀填充指定行百分比
\end{itemize}

\paragraph{\cmd{decap remove}}
默认移除峰值电流 $\le 30\,\mu$A 的无效 decap。\opt{-imax_less_than}、\opt{-dvmax_less_than} 自定义阈值。\opt{-user_decap} 仅移除用户放置 decap。\opt{-all} 移除区域内全部 decap。

\paragraph{\cmd{decap qor}}
评估 FAO 插入 decap 的 QoR：\opt{-i} 输入 instance 列表；\opt{-o} 输出目录。关联 \gsr{decap_fill}。

\paragraph{\cmd{decap drc}}
DRC 报告；\opt{-fix} 修复用户 decap 的 pin 与信号/屏蔽布线 DRC；\opt{-o <file>} 写出报告。

\subsection{示例 A--F 关键日志字段说明}
\subsubsection*{Log Field Reference for Examples A--F}

公共日志字段：
\begin{itemize}
  \item \texttt{ECO File}：输出 ECO 路径
  \item \texttt{Region}：\gsr{fao_region} 或全芯片边界
  \item \texttt{FAO Mode}：accurate / turbo
  \item \texttt{Simulation Mode}：static / dynamic
  \item \texttt{Noise Reduction}：目标压降变化百分比（负值表示放宽）
  \item \texttt{Initial Worst Static Noise}：优化前最坏静态噪声（mV）
  \item \texttt{Voltage Noise Constraints}：目标约束（mV）
  \item \texttt{Metal Usage Change Report}：各层金属面积变化百分比
  \item \texttt{DEF file for new via models}：匹配导线变更的新 via DEF
\end{itemize}

示例 C/D 另含 \texttt{Hot-spots Number}、\texttt{Fix-Window Size}、\texttt{Identified N hot-spots: (x,y)}、\texttt{Fixing Window} 坐标。示例 E/F 使用 \opt{-taper} 时修改仅限 \gsr{fao_region} 或各热点 \gsr{fix_window} 内，不横跨整片宽/高。

\subsection{手动网格修改命令速查}
\subsubsection*{Quick Reference for Manual Grid Commands}

\begin{itemize}
  \item Edit $>$ Add/Delete Pad；TCL：\cmd{fao add pad}（\opt{-window}、\opt{-metal}、\opt{-net}、\opt{-pitchX/-pitchY}、\opt{-r/-l/-c}）
  \item Edit $\rightarrow$ Add/Delete Via；叠层 via 显示为蓝色框
  \item Edit $\rightarrow$ Add Power Strap / Edit Power Strap；Include stackvias 控制相邻层或全层叠孔
  \item Edit $>$ Add Decap Cell；须 GSR \gsr{DECAP_CELL} 与 LEF 定义；APL 后从 \texttt{.cdev} 取 ESC/ESR
  \item 新增金属层：.lef LAYER、.tech metal/via、.def VIAS 段，再按常规定义 pad/strap
  \item Edit $\rightarrow$ Undo/Redo 或 \cmd{fao undo}/\cmd{fao redo}
\end{itemize}

\begin{seeAlsoBox}
第~\ref{chap:static}~章「修复 IR Drop 问题的示例流程」；附录 C「Grid Fixing and Optimization Keywords」；附录 D \cmd{export eco}/\cmd{import eco}。
\end{seeAlsoBox}

\paragraph{FAO 多 Vdd 提醒}
多 Vdd 设计中 \gsr{FAO_NETS} 每次仅处理一个电源网；须依次 \cmd{gsr set fao_nets VDD}、\cmd{gsr set fao_nets VDD_B} 分别运行 FAO，否则仅第一个网生效并产生警告。

\paragraph{mesh snscalc 使用建议}
在任意 \cmd{mesh optimize} 或 \cmd{mesh fix} 之前，建议先运行 \cmd{mesh snscalc} 查看各层灵敏度（mV/$\mu$m），优先加宽高灵敏度层导线以最有效降低最坏压降。

